\documentclass[aspectratio=169,11pt]{beamer}

\usepackage[T1]{fontenc}
\usepackage[utf8]{inputenc}
\usepackage{lmodern}
\usepackage{microtype}
\usepackage{amsmath,amssymb,mathtools,bm}
\usepackage{booktabs,array}
\usepackage{xcolor}

\definecolor{TitleBlue}{HTML}{17365D}
\definecolor{AccentTeal}{HTML}{117C88}
\definecolor{BodyGray}{HTML}{2F3A45}
\definecolor{SoftGray}{HTML}{F2F5F7}
\definecolor{SoftBlue}{HTML}{EAF1F8}
\definecolor{SoftTeal}{HTML}{E8F5F5}
\definecolor{WarningRed}{HTML}{9B2C2C}

\setbeamercolor{normal text}{fg=BodyGray,bg=white}
\setbeamercolor{frametitle}{fg=TitleBlue,bg=white}
\setbeamercolor{block title}{fg=white,bg=TitleBlue}
\setbeamercolor{block body}{fg=BodyGray,bg=SoftBlue}
\setbeamercolor{block title alerted}{fg=white,bg=WarningRed}
\setbeamercolor{block body alerted}{fg=BodyGray,bg=SoftGray}
\setbeamercolor{block title example}{fg=white,bg=AccentTeal}
\setbeamercolor{block body example}{fg=BodyGray,bg=SoftTeal}

% The compact title size is intentional: these are dense personal study slides.
\setbeamerfont{normal text}{size=\fontsize{12.2}{14.6}\selectfont}
\setbeamerfont{frametitle}{size=\fontsize{12.2}{14.6}\selectfont,series=\bfseries}
\setbeamerfont{block title}{size=\fontsize{11.4}{13.5}\selectfont,series=\bfseries}
\setbeamerfont{block body}{size=\fontsize{10.9}{13.0}\selectfont}
\setbeamerfont{footline}{size=\fontsize{8.6}{10.0}\selectfont}

\setbeamertemplate{navigation symbols}{}
\setbeamertemplate{itemize item}{\color{AccentTeal}$\blacktriangleright$}
\setbeamertemplate{itemize subitem}{\color{AccentTeal}$\bullet$}
\setlength{\leftmargini}{1.05em}
\setlength{\leftmarginii}{1.0em}
\setlength{\parskip}{0.03em}

\newcommand{\bracket}[2]{\left[#1,#2\right]}
\newcommand{\dd}{\,\mathrm{d}}
\newcommand{\avg}[1]{\left\langle #1\right\rangle}
\newcommand{\Rpsi}{R_{\psi}}
\newcommand{\Romega}{R_{\omega}}
\newcommand{\vect}[1]{\bm{#1}}
\newcommand{\OmegaD}{\Omega=[0,2\pi)\times[0,2\pi)}
\newcommand{\stateop}{(\psi_0,U_0,\eta,\nu,t)\mapsto(\psi_t,U_t)}
\newcommand{\currentmodule}{Module 1}
\newcommand{\moduleprogress}{1 of 4}
\newcommand{\setmodule}[2]{%
  \renewcommand{\currentmodule}{#1}%
  \renewcommand{\moduleprogress}{#2}%
}

\setbeamertemplate{frametitle}{%
  \vspace{0.06cm}
  \usebeamerfont{frametitle}\insertframetitle\par
  \vspace{-0.20cm}
  {\color{AccentTeal}\rule{\textwidth}{0.55pt}}
}

\setbeamertemplate{footline}{%
  \leavevmode
  \hbox{%
    \begin{beamercolorbox}[wd=.48\paperwidth,ht=2.6ex,dp=1.05ex,leftskip=.45cm]{author in head/foot}%
      \usebeamerfont{footline}\color{TitleBlue}\textbf{PINN\_M3DC1}\hspace{0.45em}|\hspace{0.45em}\currentmodule
    \end{beamercolorbox}%
    \begin{beamercolorbox}[wd=.52\paperwidth,ht=2.6ex,dp=1.05ex,rightskip=.45cm plus1fil]{date in head/foot}%
      \usebeamerfont{footline}\color{BodyGray}\moduleprogress\hspace{0.8em}\textbullet\hspace{0.8em}Slide \insertframenumber\ of \inserttotalframenumber
    \end{beamercolorbox}%
  }
}

\setbeamersize{text margin left=0.60cm,text margin right=0.60cm}
\AtBeginDocument{%
  \setlength{\abovedisplayskip}{0.22em}%
  \setlength{\belowdisplayskip}{0.22em}%
  \setlength{\abovedisplayshortskip}{0.16em}%
  \setlength{\belowdisplayshortskip}{0.16em}%
}

\begin{document}

% --------------------------------------------------------------------------
% Module 1: four slides
% --------------------------------------------------------------------------

\setmodule{Module 1}{1 of 4}
\begin{frame}[t]{Module 1: reduced resistive-MHD state}
\vspace{-0.15em}
The project physics is a strict two-dimensional, incompressible, doubly periodic reduced-MHD model:
\[
\OmegaD,\qquad t\in[0,T],\qquad \partial_z=0.
\]
The primary scalar fields are magnetic flux and stream function,
\[
\vect q(x,y,t)=(\psi,U),
\]
with derived quantities
\[
\vect B_\perp=\nabla\psi\times\widehat{\vect z}=(\psi_y,-\psi_x),\qquad
\vect v_\perp=\widehat{\vect z}\times\nabla U=(-U_y,U_x),
\]
\[
J=-\nabla^2\psi,\qquad \omega=\nabla^2U,
\qquad \nabla\cdot\vect B_\perp=\nabla\cdot\vect v_\perp=0.
\]
The surrogate receives a complete initial field pair and predicts the later field pair.
\end{frame}

\setmodule{Module 1}{2 of 4}
\begin{frame}[t]{Module 1: brackets and governing equations}
\vspace{-0.15em}
The two-dimensional bracket is
\[
\bracket{f}{g}=f_xg_y-f_yg_x,
\qquad
\vect v_\perp\cdot\nabla f=\bracket{U}{f}.
\]
Thus nonlinear advection appears compactly as a Poisson bracket. The frozen sign convention is
\[
\boxed{\psi_t+\bracket{U}{\psi}=\eta\nabla^2\psi,}
\]
\[
\boxed{\omega_t+\bracket{U}{\omega}=\bracket{J}{\psi}+\nu\nabla^2\omega,}
\qquad
J=-\nabla^2\psi,
\quad
\omega=\nabla^2U.
\]
Interpretation:
\[
\bracket{U}{\psi}:\text{ flux advection},\quad
\eta\nabla^2\psi:\text{ resistive diffusion},
\]
\[
\bracket{J}{\psi}:\text{ magnetic forcing},\quad
\nu\nabla^2\omega:\text{ viscous vorticity diffusion}.
\]
\end{frame}

\setmodule{Module 1}{3 of 4}
\begin{frame}[t]{Module 1: scales, gauges, and energy}
\vspace{-0.15em}
Alfven normalization gives
\[
V_A=\frac{B_0}{\sqrt{\mu_0\rho_0}},
\qquad
t_A=\frac{L_0}{V_A},
\qquad
\eta=\frac{\eta_m}{L_0V_A},
\qquad
\nu=\frac{\nu_d}{L_0V_A}.
\]
Dimensionless control numbers are
\[
S=\eta^{-1},\qquad \mathrm{Re}=\nu^{-1},\qquad \mathrm{Pm}=\frac{\nu}{\eta}.
\]
Periodic gauges and means:
\[
\avg{U}(t)=0,
\qquad
\frac{\mathrm d}{\mathrm dt}\avg{\psi}=0,
\qquad
\avg J=\avg\omega=0.
\]
The verification energy law is
\[
\boxed{\frac{\mathrm d}{\mathrm dt}\frac12\int_\Omega(|\nabla\psi|^2+|\nabla U|^2)\dd A
=-\eta\int_\Omega J^2\dd A-\nu\int_\Omega\omega^2\dd A\le0.}
\]
\end{frame}

\setmodule{Module 1}{4 of 4}
\begin{frame}[t]{Module 1: prediction task and claim boundary}
\vspace{-0.15em}
The learned map is
\[
\boxed{\stateop.}
\]
The inputs have direct physical roles:
\[
(\psi_0,U_0):\text{ initial magnetic and flow state},
\qquad
(\eta,\nu):\text{ diffusion strengths},
\]
\[
t:\text{ requested prediction time}.
\]
The outputs are complete numerical fields:
\[
\widehat\psi(x,y,t),\qquad \widehat U(x,y,t).
\]
The physical model does not change when training changes from data-only to physics-informed.

\begin{alertblock}{Claim boundary}
The project demonstrates a surrogate workflow for an independent two-field reduced-MHD model. It does not yet predict genuine M3D-C1 output.
\end{alertblock}
\end{frame}

% --------------------------------------------------------------------------
% Module 2: two slides
% --------------------------------------------------------------------------

\setmodule{Module 2}{1 of 2}
\begin{frame}[t]{Module 2: one simulation creates one data case}
\vspace{-0.15em}
Choose the physical settings and initial fields:
\[
(\psi_0,U_0,\eta,\nu).
\]
The solver then uses the reduced-MHD equations from Module 1 to calculate
\[
(\psi,U,J,\omega)\quad\text{at each saved time.}
\]
One case file stores
\[
x, y, t, \psi, U, J, \omega, \eta, \nu, A_0,
\]
plus the initial conditions and solver settings.

\begin{block}{Core idea}
The saved numerical arrays are the training data. Plots are only pictures used to inspect those arrays.
\end{block}
\end{frame}

\setmodule{Module 2}{2 of 2}
\begin{frame}[t]{Module 2: keep complete simulations together}
\vspace{-0.15em}
Each saved time can become one learning example:
\[
\boxed{
(\psi_0,U_0,\eta,\nu,t_k)
\longrightarrow
(\psi_k,U_k).
}
\]
All times from one simulation stay together in training, validation, or testing. This prevents the same physical evolution from appearing on both sides of the comparison.

Before use, each case must pass simple-solution, grid, time-step, periodic-boundary, energy, and save/reload checks.

\begin{alertblock}{Claim boundary}
These are reduced-MHD simulation data, not M3D-C1 results.
\end{alertblock}
\end{frame}

% --------------------------------------------------------------------------
% Module 3: five slides
% --------------------------------------------------------------------------

\setmodule{Module 3}{1 of 5}
\begin{frame}[t]{Module 3: one model contains the baseline and physics-informed versions}
\vspace{-0.15em}
The neural network predicts
\[
\mathcal G_\theta:(\psi_0,U_0,\eta,\nu,t)
\longmapsto(\widehat\psi,\widehat U).
\]
The data loss compares these predictions with Module 2:
\[
\boxed{
\mathcal L_{\rm data}
=\avg{(\widehat\psi-\psi)^2}
+\avg{(\widehat U-U)^2}.
}
\]
The baseline is simply
\[
\mathcal L_{\rm total}=w_d\mathcal L_{\rm data}.
\]

\begin{block}{Why the separate baseline module disappeared}
The physics-informed version uses the same inputs, outputs, and model. It only turns on additional loss terms.
\end{block}
\end{frame}

\setmodule{Module 3}{2 of 5}
\begin{frame}[t]{Module 3: magnetic-flux residual}
\vspace{-0.15em}
Move the predicted magnetic-flux equation to the left:
\[
\boxed{
\Rpsi
=\widehat\psi_t
+\bracket{\widehat U}{\widehat\psi}
-\eta\nabla^2\widehat\psi.
}
\]
The target is $\Rpsi=0$ at every location and time:
\[
\boxed{
\mathcal L_{{\rm PDE},\psi}
=\avg{\left(\frac{\Rpsi}{S_\psi}\right)^2}.
}
\]
Physical balance:
\[
\underbrace{\widehat\psi_t}_{\text{time change}}
+\underbrace{\bracket{\widehat U}{\widehat\psi}}_{\text{flow transports flux}}
-\underbrace{\eta\nabla^2\widehat\psi}_{\text{resistive diffusion}}
=0.
\]
This term asks for more than a similar-looking field: its time change, advection, and diffusion must balance.
\end{frame}

\setmodule{Module 3}{3 of 5}
\begin{frame}[t]{Module 3: vorticity residual}
\vspace{-0.15em}
First derive current and vorticity from the predicted fields:
\[
\widehat J=-\nabla^2\widehat\psi,
\qquad
\widehat\omega=\nabla^2\widehat U.
\]
Then form
\[
\boxed{
\Romega
=\widehat\omega_t
+\bracket{\widehat U}{\widehat\omega}
-\bracket{\widehat J}{\widehat\psi}
-\nu\nabla^2\widehat\omega.
}
\]
and penalize
\[
\boxed{
\mathcal L_{{\rm PDE},\omega}
=\avg{\left(\frac{\Romega}{S_\omega}\right)^2}.
}
\]
This enforces balance among vorticity change, flow advection, magnetic forcing, and viscous diffusion. The two PDE losses remain separate so that one equation cannot hide the other.
\end{frame}

\setmodule{Module 3}{4 of 5}
\begin{frame}[t]{Module 3: initial, boundary, and gauge constraints}
\vspace{-0.15em}
The initial state selects the correct solution of the PDE:
\[
\mathcal L_{\rm IC}
=\avg{(\widehat\psi(0)-\psi_0)^2}
+\avg{(\widehat U(0)-U_0)^2}.
\]
The stream-function gauge removes an arbitrary constant:
\[
\mathcal L_{\rm gauge}=\avg{\widehat U}^{,2}.
\]
For a periodic Fourier representation,
\[
\boxed{\mathcal L_{\rm BC}=0}
\]
because periodicity is built in. Two other constraints are also automatic:
\[
\nabla\cdot\widehat{\vect B}_\perp=0,
\qquad
\nabla\cdot\widehat{\vect v}_\perp=0,
\]
because the fields are constructed from $\widehat\psi$ and $\widehat U$.
\end{frame}

\setmodule{Module 3}{5 of 5}
\begin{frame}[t]{Module 3: final proof-of-concept objective}
\vspace{-0.15em}
The recommended first physics-informed loss is
\[
\boxed{
\mathcal L_{\rm total}
=w_d\mathcal L_{\rm data}
+w_\psi\mathcal L_{{\rm PDE},\psi}
+w_\omega\mathcal L_{{\rm PDE},\omega}
+w_0\mathcal L_{\rm IC}
+w_g\mathcal L_{\rm gauge}.
}
\]
Optional tests can add
\[
\mathcal L_{\rm mean},\qquad
\mathcal L_{\rm energy},\qquad
\mathcal L_{J,\omega}.
\]
The weights prevent one numerically large term from overwhelming the others.

\begin{exampleblock}{What this presentation establishes}
It shows exactly how simulation agreement, the two MHD equations, and known conditions are combined to construct a physics-informed surrogate. Accuracy or speed improvements require later trained results.
\end{exampleblock}
\end{frame}

\end{document}
